\chapter[Conferenz-Anzeigen]{Conferenz-Anzeigen}
\quellseite{12}
Die Post Dak.-Conferenz tagt, D. V., am 12. und 13. Januar bei P. W. C. Dube zu Hochkirch (P. O. Rock), Tex. Arbeiten: XIV. Artikel der Augsburgischen Confession; P. Sieck. Predigerlehren: P. Greßens. Exegese über das 1. Capitel des Galaterbriefes: P. Bernthal. Ein Confirmations-Examen: P. Birkmann. Prediger: P. Kramer (P. Tritsche). Beichtredner: P. Tegeler (P. Heimeier). An-, resp. Abmeldung nebst Angabe, ob man in Thorndale oder in Taylor abzusteigen gedenkt, erbeten. W. C. Dube, Secr.

Die Quincy-Specialconferenz versammelt sich, will’s Gott, vom 19. bis 21. Januar in Warsaw, Ill. P. Rabe bittet um zeitige An- oder Abmeldung. W. Schaller.

Die Saginaw Valley-Pastoralconferenz tagt am 26. und 27. Januar bei Mount Pleasant, Mich. Arbeiten: Mayer, Umbach, Boß, Grüber. Prediger: Umbach (Jüngst). An- und Abmeldung erbeten bis zum 22. Januar. H. Grüber.

Die Süd-Michigan-Pastoralconferenz versammelt sich, D. V., am 26. und 27. Januar in der Gemeinde P. E. C. Bauers zu Wyanotte, Mich. Arbeiten: 1. Sermonet: P. Spiegle; 2. Exegese über 2 Tim. 2: P. Röder; 3. Exegese über 2 Tim. 3: P. Gendick; 4. Predigten zur Kritik: P. Behrens über das Evangelium des 3. Sonnt. d. Adv. und P. Claus jun. über das Evangelium des 4. Sonnt. d. Adv.; 5. Katechese: Vom Nutzen des heiligen Abendmahls, Frage 335—337 des Synodal-Katechismus: P. Tews; 6. Prediger: P. Scheler (P. Riedel); 7. Beichtredner: P. Lübbe (P. Hagen). Rechtzeitige Anmeldung ist erwünscht. B. Richter, Secr.
